\section{Vortex sector, competing infrared scales, and \BKT{} implications}
\label{sec:vortex_rg}

The excitonic order parameter admits a polar decomposition before any renormalization-group analysis,
\begin{equation}
\Psi(\mathbf{r})=\rho(\mathbf{r})\,e^{i\theta(\mathbf{r})}.
\label{eq:polar}
\end{equation}
The locked gap map of Section~\ref{sec:gap} lives in the amplitude $\rho$. The vortex sector of the theory lives in the phase $\theta$. Together the two sectors furnish a competing-scale structure: a saddle-node fold in $\rho$, and a Kosterlitz--Thouless unbinding surface in the stiffness of $\theta$. When a physical drive reaches the KT surface first, the implied gap closing is a Berezinskii--Kosterlitz--Thouless essential singularity. That implication is part of the theory; it is stated as hypothesis H-VRG1 and is evaluated below by the computational suite C1--C6.

Option~A* freezes the toroidal cavity, so the vacuum amplitude---and therefore $\alpha$---does not run with the KT scale $\ell=\ln(L/a_c)$. The suite therefore scans frozen $\alpha$ along the locked physical branch $\delta_+(\alpha)$. Kosterlitz's coupling is $K=\kappa/T_{\mathrm{eff}}$, with critical value $K_c=2/\pi$. No bare transition temperature is inserted by hand, and no chain-rule $\beta_\alpha$ is integrated.

\begin{center}
\fcolorbox{oiBlue}{oiBlue!6}{%
\parbox{0.94\linewidth}{\small
\textbf{Defect vocabulary.}
A \textbf{hopfion} is a map $S^3\to S^2$ with fractional Hopf charge $Q_H=1/N$, realized here as a toroidal Casimir cavity.
A \textbf{2D phase vortex} is a point defect in the moir\'e plane with $\oint\nabla\theta\cdot d\mathbf{l}=2\pi n$; this is the \BKT{} defect of the present section.
A \textbf{Chern--Simons fluxon} is a $2\pi$ winding that binds flux $\Phi_{\mathrm{CS}}=2\pi/N$ at level $k=2N$.
The three objects are distinct.
}}
\end{center}

\subsection{Two-sector decomposition}
\label{sec:two_sector}

Substituting Eq.~\eqref{eq:polar} into the bulk \TQGL{} energy of Section~\ref{sec:tqgl} separates a Higgs potential for $\rho$ from a phase stiffness for $\theta$,
\begin{equation}
\mathcal{F}_{\mathrm{bulk}}[\rho,\theta]
=\int\mathrm{d}^2r\Biggl[
\alpha_{\mathrm{GL}}\rho^2+\frac{\beta}{2}\rho^4
+\kappa_{\mathrm{GL}}\bigl(|\nabla\rho|^2+\rho^2|\nabla\theta|^2\bigr)
\Biggr],
\label{eq:two_sector}
\end{equation}
up to the Chern--Simons, Berry, and Skyrme terms retained in the parent functional. Here $\kappa_{\mathrm{GL}}$ is the Ginzburg--Landau gradient weight written $\kappa$ in Section~\ref{sec:tqgl}, and $\alpha_{\mathrm{GL}}$ is the Landau mass, distinct from the gap-map coupling $\alpha$ of Section~\ref{sec:gap}. Equation~\eqref{eq:two_sector} captures the algebra of the split. The fold of the amplitude map at $\alpha_c=\delta_c=\alphac$, $\Delta_c=\Deltac$ is a locked input. Vortex unbinding is the content of the phase sector.

The live-gap estimator $\Delta_{\mathrm{live}}[\psi]=\Delta_0\sqrt{n_w[\psi]/n_{\mathrm{ref}}}$ of Section~\ref{sec:dynamics} is an amplitude proxy. A field whose modulus remains finite while $\theta$ disorders still returns a nonzero $\Delta_{\mathrm{live}}$. Amplitude and phase diagnostics therefore complement one another: a $1$\,ps trajectory with $\Delta_{\mathrm{live}}(1\,\mathrm{ps})=\Deltaliveone$ constrains $\rho$, while $g(r)$ and the stiffness flow constrain $\theta$.

\begin{figure}[htbp]
\centering
% Two-sector split: amplitude fold vs 2D phase vortex RG
\begin{tikzpicture}[
  font=\sffamily\footnotesize,
  box/.style={draw, rounded corners=2pt, align=center, inner sep=6pt,
    text width=3.55cm, minimum height=1.35cm, line width=0.8pt},
  amp/.style={box, draw=oiBlue, fill=oiBlue!10},
  ph/.style={box, draw=oiOrange, fill=oiOrange!12},
  lock/.style={box, draw=oiGreen, fill=oiGreen!10, text width=3.2cm},
  conj/.style={box, draw=oiVerm, fill=oiVerm!10, dashed, text width=3.2cm},
]
\node[font=\sffamily\small\bfseries] at (0,3.55) {Amplitude $\rho$};
\node[font=\sffamily\small\bfseries] at (7.2,3.55) {Phase $\theta$};
\node[amp] (psi) at (3.6,4.85) {$\Psi=\rho\,e^{i\theta}$\\[-1pt]{\scriptsize before RG}};
\node[amp] (fold) at (0,2.15) {gap map $\delta=1+\alpha(\ln\delta+\ell)$};
\node[lock] (sn) at (0,0.45) {\statuslocked{} saddle-node\\$\alpha_c=\delta_c=\alphac$};
\node[ph] (xy) at (7.2,2.15) {stiffness $\kappa$, fugacity $y$\\Villain / Coulomb gas};
\node[conj] (kt) at (7.2,0.45) {\statusconjecture{} H-VRG1\\KT surface $K=2/\pi$};
\draw[-{Stealth}, oiBlue, thick] (psi.south) -- (fold.north);
\draw[-{Stealth}, oiOrange, thick] (psi.south) -- (xy.north);
\draw[-{Stealth}, oiBlue] (fold.south) -- (sn.north);
\draw[-{Stealth}, oiOrange] (xy.south) -- (kt.north);
\node[font=\scriptsize, text width=10.5cm, align=center, oiVerm!70!black]
  at (3.6,-1.05) {Live-gap estimator $\Delta_0\sqrt{n_w/n_{\mathrm{ref}}}$ reads $\rho$, not vortex unbinding.
  Chern--Simons is \statustruncated{} in the GPE and in the default XY--KT scan.};
\end{tikzpicture}

\caption{\textbf{Two-sector split.} The locked fold is an amplitude phenomenon. Vortex RG is a phase phenomenon. The live-gap estimator reads $\rho$. Chern--Simons is truncated from the default XY--KT scan, matching the GPE truncation of Section~\ref{sec:dynamics}.}
\label{fig:two_sector}
\end{figure}

\subsection{C1: phase stiffness}
\label{sec:stiffness}

The 2D phase stiffness (superfluid density) extracted from Eq.~\eqref{eq:two_sector} is not $\kappa_{\mathrm{GL}}$ itself. Matching $\kappa_{\mathrm{GL}}|\nabla\Psi|^2$ to the conventional $(\rho_s/2)|\nabla\theta|^2$ with $\Psi=\sqrt{n_0}\,\rho\,e^{i\theta}$ gives
\begin{equation}
\kappa\equiv\rho_s=2 n_0\kappa_{\mathrm{GL}}\langle\rho^2\rangle+\kappa_{\mathrm{geom}},
\label{eq:kappa}
\end{equation}
where $\kappa_{\mathrm{GL}}=\Delta_0 a_H^2$ follows from the truncated GPE kinetic match and $\langle\rho^2\rangle=\delta_+(\alpha)^2$ is read from the locked physical branch through the live-gap identification. The density unit is the healing-disk ansatz $n_0=1/(\pi a_H^2)$, which converts dimensionless $|\psi|^2\sim 1$ into a 2D density and makes the conventional piece at unit amplitude equal to $2\Delta_0/\pi$, independent of $a_H$. The geometric weight $\kappa_{\mathrm{geom}}$ is a frozen Peotta--T\"orm\"a floor~\cite{PeottaTorma2015},
\begin{equation}
\kappa_{\mathrm{geom}}=\frac{\Delta_0}{2\pi},
\label{eq:kappa_geom}
\end{equation}
treated as an ansatz: it is not a Brillouin-zone integral of the TEVC quantum metric, and it need not vanish with $\delta$. From here on in this section, $\kappa$ means the stiffness~\eqref{eq:kappa}. A uniform phase-twist check on the same conventional piece agrees with $2 n_0\kappa_{\mathrm{GL}}\langle\rho^2\rangle$ to machine precision.

On this primary vortex-RG closure---healing-disk $n_0$ plus the frozen geometric floor~\eqref{eq:kappa_geom}---one finds $\kappa_{\mathrm{UV}}\approx\kappauvfloor$ at $\alpha=0$ and $\kappa(\alpha_c)\approx 0.108$\,meV at the fold. With $T_{\mathrm{eff}}=T^*=\Tstar$ the Kosterlitz couplings are $K(\alpha=0,T^*)\approx\Kuvfloor$ and $K(\alpha_c,T^*)\approx\Kfoldfloor$, both far above $K_c=2/\pi\approx 0.637$. A Qi--Wu--Zhang toy metric (Section~\ref{sec:c1_trackb}) is a second labeled C1 estimator ($\kappa_{\mathrm{UV}}^{\mathrm{toy}}=\kappauvtoy$, $K_{\mathrm{fold}}(T^*)=\KfoldTstar$); both miss KT at $T^*$ before the fold. A conventional-only variant with $\kappa_{\mathrm{geom}}=0$ is retained as a labeled alternative; its ultraviolet stiffness is $2\Delta_0/\pi\approx 0.382$\,meV. The energy of an isolated 2D phase vortex in a system of size $L$ with core cutoff $a_c$ is
\begin{equation}
E_v=\pi\kappa\ln(L/a_c)+E_{\mathrm{core}}.
\label{eq:Ev}
\end{equation}
The gap $\Delta$ sets a scale for $\rho$ and for $a_c\sim\hbar c_s/\Delta$. \BKT{} physics is a statement about whether $\kappa/T_{\mathrm{eff}}$ flows through the universal jump
\begin{equation}
\kappa(T_{\mathrm{BKT}}^-)=\frac{2T_{\mathrm{BKT}}}{\pi}.
\label{eq:universal_jump}
\end{equation}
A finite $\Delta_c$ at the fold is compatible with a still-superfluid phase, with a plasma of 2D phase vortices, or with a rounded quantum crossover; the amplitude map and the stiffness flow decide jointly.

\subsection{Coulomb gas and Villain action}
\label{sec:villain}

The effective theory for 2D phase vortices is the XY/Villain model~\cite{Kosterlitz1973,Villain1975,Jose1977}, equivalently the integer Coulomb gas
\begin{equation}
\frac{H}{T_{\mathrm{eff}}}
=2\pi K\sum_{i<j}n_i n_j\ln\frac{|r_i-r_j|}{a_c}
+\frac{E_{\mathrm{core}}}{T_{\mathrm{eff}}}\sum_i n_i^2,
\label{eq:villain}
\end{equation}
with spin-wave piece $S_{\mathrm{sw}}[\theta]=\frac{K}{2}\int(\nabla\theta)^2$. We adopt the Kosterlitz normalization $K=\kappa/T_{\mathrm{eff}}$ (not the Jos\'e--Kadanoff $\tilde K=\pi\kappa/T_{\mathrm{eff}}$), so that the one-loop flow
\begin{equation}
\frac{\mathrm{d}K^{-1}}{\mathrm{d}\ell}=4\pi^3 y^2,\qquad
\frac{\mathrm{d}y}{\mathrm{d}\ell}=(2-\pi K)y
\label{eq:kt_flow}
\end{equation}
has critical point $K_c=2/\pi$. Equivalently, $\pi\kappa/T_{\mathrm{eff}}=\pi K$ equals $2$ on the universal jump~\eqref{eq:universal_jump}. Equation~\eqref{eq:kt_flow} is the vortex-RG sector of the theory. Its derivation from the microscopic \TQGL{} functional proceeds through a Villain reduction, which we state as an ansatz: the phase sector is governed by~\eqref{eq:villain}--\eqref{eq:kt_flow} once amplitude fluctuations have been integrated at the scale $a_c$. The flow equations are derived as an effective theory; the Villain reduction from the full \TQGL{} functional is treated as an ansatz. The fugacity is the textbook form $y=\exp(-E_{\mathrm{core}}/T_{\mathrm{eff}})$.

\subsection{C2: winding-one core}
\label{sec:c2_core}

The defect that enters~\eqref{eq:Ev} is a winding-$1$ radial Gross--Pitaevskii vortex in the moir\'e plane, not a hopfion and not a Chern--Simons fluxon. Integrating the truncated GP energy on a disk and subtracting the infrared logarithm at the healing-length cutoff $a_c=\xi=8$\,nm yields $E_{\mathrm{core}}/\kappa\approx 0.343$, with the cutoff itself an ansatz. That ratio sets the bare fugacity along the branch, $y=\exp\bigl[-(E_{\mathrm{core}}/\kappa)K\bigr]$. Because $E_{\mathrm{core}}/\kappa$ is $\mathcal{O}(1)$ rather than $\mathcal{O}(10)$, finite $y$ can matter even when the bare coupling still lies well above $2/\pi$: the infrared plasma is a separatrix phenomenon, not merely a $K<K_c$ crossing of the ultraviolet data.

\subsection{Competing infrared scales}
\label{sec:competing}

Two surfaces compete in the $(\alpha,K,y)$ space of a driven TEVC.

\begin{theorem}[Competing-scale criterion]
\label{thm:competing}
On the locked amplitude theory the fold surface is
\begin{equation}
\alpha=\alpha_c,\qquad \Delta=\Delta_c>0.
\label{eq:fold_surface}
\end{equation}
On the Villain target~\eqref{eq:kt_flow} the KT surface is
\begin{equation}
K=2/\pi,\qquad y\to 0,
\label{eq:kt_surface}
\end{equation}
with $\Delta\to 0$ only if $\Delta$ tracks a positive power of $\kappa$ after unbinding. Hypothesis H-VRG1 holds if and only if the physical drive path hits the KT surface \emph{before} the fold. If $\kappa(A)$ remains larger than $\kappa_{\mathrm{BKT}}$ at $\alpha=\alpha_c$ and finite-$y$ flow stays bound, the fold wins and the gap map retains a finite $\Delta_c$.
\end{theorem}

Whether a given protocol satisfies it is a property of $\kappa(\alpha)$, $E_{\mathrm{core}}/\kappa$, and $T_{\mathrm{eff}}$.

\begin{hypothesis}[H-VRG1]
Including a 2D phase-vortex fugacity in the RG of $\theta=\arg\Psi$ converts the locked saddle-node fold into a \textup{\BKT{}} essential singularity $\Delta\sim\exp(-C/\sqrt{A_c-A})$ with continuous gap vanishing, whenever the drive path of Theorem~\ref{thm:competing} reaches the KT surface first.
\end{hypothesis}

H-VRG1 is a conjecture. The C1 stiffness is a computed closure with ansatz inputs, not a Brillouin-zone theorem, so the competing-scale diagram below is the implication of that closure.

\subsection{Chern--Simons statistics: truncation or anyon plasma}
\label{sec:anyon}

The parent functional contains a Chern--Simons term at level $k=2N$, which attaches flux $\Phi_{\mathrm{CS}}=2\pi/N$ to a $2\pi$ winding. Ordinary XY--KT assumes bosonic integer 2D phase vortices. An anyonic plasma changes the Coulomb coupling.

Two consistent options exist. The first, adopted for the default scan and matching the GPE of Section~\ref{sec:dynamics}, is to \emph{truncate} Chern--Simons, Berry, and Skyrme from the vortex sector: the flow~\eqref{eq:kt_flow} then applies to ordinary 2D phase vortices, and CS fluxons are outside the calculation. The second is to restore level $k=2N$ and treat 2D phase vortices as anyons with statistical angle
\begin{equation}
\theta_{\mathrm{stat}}=\frac{\pi}{N}\qquad(N=5\Rightarrow \theta_{\mathrm{stat}}=\pi/5).
\label{eq:theta_stat}
\end{equation}
A Wen--Zee-inspired shift of the KT criterion~\cite{WenZee1992,Wen1995} then reads
\begin{equation}
K_c^{\mathrm{anyon}}=\frac{2}{\pi}\Bigl(1-\frac{\theta_{\mathrm{stat}}}{\pi}\Bigr)^2
=\frac{2}{\pi}\Bigl(1-\frac{1}{N}\Bigr)^2,
\label{eq:Kc_anyon}
\end{equation}
which for $N=5$ lowers $K_c$ from $2/\pi\approx 0.637$ to $\approx 0.407$. A lower $K_c$ moves the KT surface toward weaker stiffness. Equation~\eqref{eq:Kc_anyon} remains conjectural unless the anyon-Coulomb-gas algebra is derived from the \TQGL{} Chern--Simons term; it is displayed so that the truncation is explicit, and it is never mixed into the truncated default.

\subsection{C3: competing-scale scan}
\label{sec:vrg_python}

On the primary closure---healing-disk $n_0$ plus the frozen geometric floor---at $T^*=\Tstar$, Kosterlitz flow at frozen $\alpha$ remains bound all the way to the fold. The infrared coupling at $\alpha_c$ is $K(\alpha_c)\approx 10.57\gg 2/\pi$: the fold is reached first. Restoring the labeled anyon threshold~\eqref{eq:Kc_anyon} does not change the winner on this $\kappa(\alpha)$. The same $E_{\mathrm{core}}/\kappa\approx 0.34$ that would be dangerous near $K\sim 1$ is harmless at $K\sim 10$, because $y=\exp(-0.34 K)$ is then exponentially small.

The conventional-only variant ($\kappa_{\mathrm{geom}}=0$) at the same $T^*$ is different. Bare $K$ still stays above $2/\pi$ through $\alpha_c$, so there is no ultraviolet $K<K_c$ crossing. Finite-$y$ flow, however, hits the plasma separatrix at $\alpha/\alpha_c\approx 0.945$: KT before the fold, as an infrared phenomenon. The computed core ratio is large enough for that separatrix even though the geometric floor is what had kept the primary branch bound.

Raising $T_{\mathrm{eff}}$ opens a temperature window for the vortex sector even on the primary closure. Below $\sim 50$\,mK both closures remain bound out to $\alpha_c$. At $T^*$ only conventional-only unbinds. At $T_{\mathrm{eff}}=0.5$\,K the primary branch already plasmas at $\alpha/\alpha_c\approx 0.819$; at $1$\,K the plasma station has moved in to $\alpha/\alpha_c\approx 0.302$. That window is the physically interesting content of the sector: whether a laboratory protocol sits on the fold side or the KT side of Theorem~\ref{thm:competing} is a joint question of stiffness closure and temperature, not an algebraic identity.

\begin{figure}[htbp]
\centering
% Textbook KT flow at frozen α: primary bound vs conv-only plasma (C3)
\begin{tikzpicture}
\begin{axis}[
  width=0.92\linewidth,
  height=6.4cm,
  xmode=log,
  ymode=log,
  xlabel={$K=\kappa/T_{\mathrm{eff}}$},
  ylabel={fugacity $y$},
  xmin=0.2, xmax=60,
  ymin=1e-6, ymax=2,
  legend style={draw=none, at={(0.98,0.98)}, anchor=north east, font=\small, cells={anchor=west}},
  tick label style={font=\small},
  label style={font=\small},
  clip=true,
]
\addplot[oiGreen, line width=1.5pt, mark=*, mark size=1.5pt]
  table[x index=1, y index=2, skip first n=1] {data/vortex_rg_flow_uv.dat};
\addlegendentry{primary UV (bound)}
\addplot[oiBlue, line width=1.5pt, mark=*, mark size=1.5pt]
  table[x index=1, y index=2, skip first n=1] {data/vortex_rg_flow_star.dat};
\addlegendentry{primary $\lambda^\star$ (bound)}
\addplot[oiSky, line width=1.5pt, dashed, mark=*, mark size=1.5pt, every mark/.append style={solid}]
  table[x index=1, y index=2, skip first n=1] {data/vortex_rg_flow_nearfold.dat};
\addlegendentry{primary near fold (bound)}
\addplot[oiVerm, line width=1.5pt]
  table[x index=1, y index=2, skip first n=1] {data/vortex_rg_flow_conv_plasma.dat};
\addlegendentry{conv-only IR plasma}
\draw[black!50, densely dotted, line width=0.8pt]
  (axis cs:0.63662,1e-6) -- (axis cs:0.63662,2);
\node[font=\small, oiBlue, anchor=south, rotate=90] at (axis cs:0.85,3e-4) {$K_c=2/\pi$};
\end{axis}
\end{tikzpicture}

\caption{\textbf{Kosterlitz flow at frozen $\alpha$ and $T^*$.} Primary-closure trajectories at $\alpha/\alpha_c=0.30$, at $\lambda^\star$, and near the fold all run to $y\to 0$ with $K\gg 2/\pi$ (bound). A conventional-only station at the finite-$y$ onset $\alpha/\alpha_c\approx 0.945$ runs to a plasma. Option~A* holds $\alpha$ fixed; there is no $\beta_\alpha$.}
\label{fig:vrg_flow}
\end{figure}

\begin{figure}[htbp]
\centering
% Competing IR scales along the locked physical branch (C1/C3 at T*)
\begin{tikzpicture}
\begin{axis}[
  width=0.92\linewidth,
  height=6.4cm,
  ymode=log,
  xlabel={$\alpha/\alpha_c$},
  ylabel={$K=\kappa/T^*$},
  xmin=0, xmax=1.02,
  ymin=0.3, ymax=70,
  legend style={draw=none, fill=white, fill opacity=0.85, at={(0.02,0.02)}, anchor=south west, font=\small, cells={anchor=west}},
  tick label style={font=\small},
  label style={font=\small},
  grid=both,
  major grid style={line width=0.3pt, draw=black!20},
  minor grid style={line width=0.15pt, draw=black!8},
  minor x tick num=4,
]
\addplot[oiBlue, line width=1.5pt]
  table[x index=0, y index=1, skip first n=1] {data/competing_scales.dat};
\addlegendentry{primary $K$ (C1)}
\addplot[oiOrange, line width=1.5pt, dashed]
  table[x index=0, y index=2, skip first n=1] {data/competing_scales.dat};
\addlegendentry{conventional-only}
\addplot[oiVerm, densely dashed, line width=0.8pt]
  coordinates {(0,0.63662) (1.02,0.63662)};
\addlegendentry{$K_c=2/\pi$ (bosonic)}
\addplot[oiPurple, densely dotted, line width=0.8pt]
  coordinates {(0,0.4074) (1.02,0.4074)};
\addlegendentry{$K_c^{\mathrm{anyon}}$ (labeled)}
\addplot[oiGreen, densely dashed, line width=0.8pt]
  coordinates {(0.90,0.3) (0.90,70)};
\addlegendentry{$\lambda^\star$ ($\alpha/\alpha_c=0.90$)}
\node[font=\small, oiGreen, anchor=south west, rotate=90] at (axis cs:0.91,1.0)
  {$\alpha/\alpha_c=0.90$};
\addplot[only marks, mark=*, mark size=2.4pt, oiOrange]
  coordinates {(0.945, 4.04)};
\addlegendentry{conv.\ IR plasma $\approx 0.945\,\alpha_c$}
\end{axis}
\end{tikzpicture}

\caption{\textbf{Competing-scale diagram at $T^*$.} Kosterlitz $K=\kappa/T^*$ along the locked physical branch for the primary C1 closure and the conventional-only variant. The bosonic line $K=2/\pi$ and the labeled anyon line lie far below the primary curve; the fold is reached first. Conventional-only unbinding is a finite-$y$ infrared plasma at $\alpha/\alpha_c\approx 0.945$, marked, not a bare $K<2/\pi$ crossing.}
\label{fig:competing_scales}
\end{figure}

\begin{figure}[htbp]
\centering
% Temperature window: first plasma station vs T_eff (C3)
\begin{tikzpicture}
\begin{axis}[
  width=0.92\linewidth,
  height=5.8cm,
  xmode=log,
  xlabel={$T_{\mathrm{eff}}$ (K)},
  ylabel={first plasma $\alpha/\alpha_c$},
  xmin=0.015, xmax=1.4,
  ymin=0, ymax=1.08,
  legend style={draw=none, at={(0.02,0.12)}, anchor=south west, font=\footnotesize, cells={anchor=west}},
  tick label style={font=\small},
  label style={font=\small},
]
\addplot[oiBlue, thick, mark=*, mark size=2.2pt]
  table[x index=0, y index=2, skip first n=1] {data/Teff_scan.dat};
\addlegendentry{primary (fold $=1$)}
\addplot[oiOrange, thick, dashed, mark=square*, mark size=2.2pt]
  table[x index=0, y index=4, skip first n=1] {data/Teff_scan.dat};
\addlegendentry{conventional-only}
\addplot[oiGreen, densely dashed]
  coordinates {(0.1187,0) (0.1187,1.08)};
\addlegendentry{$T^*=118.7$\,mK}
\addplot[oiVerm, densely dotted]
  coordinates {(0.015,1) (1.4,1)};
\end{axis}
\end{tikzpicture}

\caption{\textbf{Temperature window of the vortex sector.} First $\alpha/\alpha_c$ at which frozen-$\alpha$ flow reaches an infrared plasma, or $\alpha/\alpha_c=1$ if the fold is reached first. The Schwinger crossover $T^*$ is marked. Raising $T_{\mathrm{eff}}$ to $0.5$--$1$\,K lets even the primary closure unbind by finite $y$.}
\label{fig:teff_window}
\end{figure}

\subsection{C4: vacuum driving and the essential-singularity implication}
\label{sec:unbinding_mech}

On the primary $T^*$ scan the flow stays bound, so the locked finite $\Delta_c$ stands: there is no vortex-driven continuous $\Delta\to 0$ on that closure. When a protocol does pass Theorem~\ref{thm:competing}---as the conventional-only variant does at $T^*$---increasing the vacuum amplitude $A$ (equivalently $\alpha$) walks the system along $\delta_+(\alpha)\searrow\delta_c$ while $\kappa$ and $E_{\mathrm{core}}$ fall and $y$ rises. Unbinding occurs if the finite-$y$ separatrix, or $K=2/\pi$, is reached while $\alpha$ is still less than $\alpha_c$. Identifying the KT correlation length with a gap,
\begin{equation}
\Delta\sim\hbar c_s/\xi_{\mathrm{KT}},
\label{eq:delta_xi}
\end{equation}
then produces the essential singularity
\begin{equation}
\Delta(A)\sim\exp\bigl(-C/\sqrt{A_c-A}\bigr)
\label{eq:v9}
\end{equation}
whenever $\Delta$ tracks $\sqrt{\kappa}$ or an equivalent stiffness collapse. Equation~\eqref{eq:delta_xi} is a conjectural identification; Eq.~\eqref{eq:v9} is the \BKT{} implication of the vortex sector (H-VRG1). It is a physical prediction for a protocol that passes Theorem~\ref{thm:competing}, not a replacement of the locked amplitude fold.

\subsection{Finite-size phase correlations}
\label{sec:gr}

Define the 2D phase correlator
\begin{equation}
g(r)=\bigl\langle e^{i[\theta(\mathbf{r})-\theta(0)]}\bigr\rangle.
\label{eq:gr}
\end{equation}
In a Villain spin-wave treatment (vortices omitted) $g(r)\sim r^{-\eta}$ with $\eta=1/(2\pi K)$; at $K=2/\pi$ one has $\eta=1/4$. In a plasma, $g(r)$ is exponential. An effective exponent on a torus of size $L$,
\begin{equation}
\eta_{\mathrm{eff}}(L)=-\frac{\ln g(L/2)}{\ln(L/(2a))},
\label{eq:eta_eff}
\end{equation}
approaches $1/4$ at a putative KT point and rises when vortices proliferate. On the primary $T^*$ branch, $K\sim 10$--$47$ implies a spin-wave $\eta$ of order $10^{-3}$--$10^{-2}$, far below $1/4$, consistent with a still-bound phase. Figure~\ref{fig:gr} shows that family together with the algebraic KT target and an exponential plasma toy; Fig.~\ref{fig:eta_eff} shows $\eta_{\mathrm{eff}}$ along the physical branch. A modest $L=16$ Metropolis XY run remains a program item for sampling vortices rather than spin waves. 
\begin{figure}[htbp]
\centering
% Villain spin-wave g(r) at primary λ* stiffness (C3/C1)
\begin{tikzpicture}
\begin{axis}[
  name=gplot,
  width=0.92\linewidth,
  height=5.6cm,
  xmode=log,
  ymode=log,
  xlabel={$r/a$},
  ylabel={$g(r)$},
  xmin=1, xmax=64,
  ymin=1e-3, ymax=1.2,
  legend style={draw=none, at={(0.02,0.02)}, anchor=south west, font=\footnotesize, cells={anchor=west}},
  tick label style={font=\small},
  label style={font=\small},
]
\addplot[oiBlue, thick]
  table[x index=0, y index=1, skip first n=1] {data/g_r_scaling.dat};
\addlegendentry{spin-wave at $\lambda^\star$}
\addplot[oiOrange, thick, dashed]
  table[x index=0, y index=2, skip first n=1] {data/g_r_scaling.dat};
\addlegendentry{$\eta=1/4$ (KT target)}
\addplot[oiVerm, densely dotted]
  table[x index=0, y index=3, skip first n=1] {data/g_r_scaling.dat};
\addlegendentry{plasma toy}
\end{axis}
\end{tikzpicture}

\caption{\textbf{Phase correlator $g(r)$.} Villain spin-wave decay at the primary $\lambda^\star$ stiffness, the algebraic KT target $\eta=1/4$, and an exponential plasma toy. Vortices are not sampled in the spin-wave curves.}
\label{fig:gr}
\end{figure}

\begin{figure}[htbp]
\centering
% Spin-wave η along the primary T* branch
\begin{tikzpicture}
\begin{axis}[
  width=0.92\linewidth,
  height=6.0cm,
  xlabel={$\alpha/\alpha_c$},
  ylabel={$\eta=1/(2\pi K)$},
  xmin=0, xmax=1.0,
  ymin=0, ymax=0.020,
  legend style={draw=none, at={(0.02,0.98)}, anchor=north west, font=\small, cells={anchor=west}},
  tick label style={font=\small},
  label style={font=\small},
]
\addplot[oiBlue, line width=1.5pt]
  table[x index=0, y index=2, skip first n=1] {data/eta_eff.dat};
\addlegendentry{primary spin-wave $\eta$}
% η_c = 1/4 = 0.25 is far above this range; annotate with upward arrow
\draw[oiVerm, densely dashed, line width=0.8pt, -{Stealth[length=4pt]}]
  (axis cs:0.55,0.016) -- (axis cs:0.55,0.0198);
\node[font=\small, oiVerm, anchor=west] at (axis cs:0.57,0.0175)
  {$\eta_c = \tfrac{1}{4}$ (far above)};
\end{axis}
\end{tikzpicture}

\caption{\textbf{Spin-wave $\eta=1/(2\pi K(\alpha))$ along the primary $T^*$ branch.} With $K(\alpha_c)\approx 10.57$ the exponent remains far below the KT target $\eta=1/4$. The line $\eta=1/4$ is the would-be unbinding marker.}
\label{fig:eta_eff}
\end{figure}

\subsection{C5: fold versus essential singularity}
\label{sec:c5}

On the locked physical branch, $\delta-\delta_c\sim\sqrt{t_A}$ with $t_A=(\alpha_c-\alpha)/\alpha_c$ is a theorem of the amplitude map. A Gaussian comparison of $\tfrac12\ln t_A$ against $1/\sqrt{t_A}$ on $\ln(\delta-\delta_c)$ over $\gtrsim 1.5$ decades in $t_A$ returns a free slope $0.542$ (versus the theorem value $1/2$) and prefers the square root by $\Delta\ell\ell\approx +431$. That comparison is a consistency check on fold data. Vortex-implied $\Delta$ from Eq.~\eqref{eq:delta_xi} is not available on the primary $T^*$ scan, which never unbinds, so the essential-singularity curve in Fig.~\ref{fig:delta_tA} is the H-VRG1 target overlaid on a map that does not contain vortices. The preference for $\sqrt{t_A}$ confirms the locked fold; it is not a statement that a KT surface cannot be hit on a different closure or at a higher $T_{\mathrm{eff}}$.

\begin{figure}[htbp]
\centering
% Fold √ scaling vs H-VRG1 essential singularity (C5)
\begin{tikzpicture}
\begin{axis}[
  width=0.92\linewidth,
  height=6.2cm,
  xmode=log,
  ymode=log,
  xlabel={$t_A=(\alpha_c-\alpha)/\alpha_c$},
  ylabel={$\delta$},
  xmin=1e-3, xmax=1,
  ymin=1e-5, ymax=1.2,
  legend style={draw=none, fill=white, fill opacity=0.85, at={(0.98,0.98)}, anchor=north east, font=\small, cells={anchor=west}},
  tick label style={font=\small},
  label style={font=\small},
]
\addplot[oiBlue, line width=1.5pt]
  table[x index=0, y index=1, skip first n=1] {data/delta_vs_tA.dat};
\addlegendentry{physical branch (\statuslocked)}
\addplot[oiGreen, densely dashed, line width=1.5pt]
  table[x index=0, y index=2, skip first n=1] {data/delta_vs_tA.dat};
\addlegendentry{$\delta_c+B\sqrt{t_A}$}
\addplot[oiVerm, densely dotted, line width=1.5pt]
  table[x index=0, y index=3, skip first n=1] {data/delta_vs_tA.dat};
\addlegendentry{H-VRG1 $e^{-C/\sqrt{t_A}}$}
\addplot[oiOrange, densely dotted, line width=0.8pt]
  coordinates {(1e-3,0.1818) (1,0.1818)};
\end{axis}
\end{tikzpicture}

\caption{\textbf{Amplitude-map scaling on the locked branch.} Physical $\delta(t_A)$ follows $\delta_c+B\sqrt{t_A}$. The dotted essential singularity is the H-VRG1 target of the vortex sector; vortex-implied $\Delta$ is not constructed on the primary $T^*$ scan.}
\label{fig:delta_tA}
\end{figure}

\subsection{C6: thermal versus quantum unbinding}
\label{sec:thermal_quantum}

Classical KT flow requires $T>0$. The TEVC condensation window is $\lesssim 1$\,K; the Schwinger crossover is $T^*=\Tstar$; the live GPE of Section~\ref{sec:dynamics} is $T=0$. At strictly zero temperature the thermal fugacity $y=e^{-E_{\mathrm{core}}/T}$ vanishes and textbook KT does not run. Quantum phase slips, the 2D Bose--Hubbard/Villain dual~\cite{FisherWeichman1989}, or dissipative Ambegaokar--Halperin--Nelson--Siggia (AHNS) dynamics~\cite{Ambegaokar1980,HalperinNelson1978} are then the relevant mechanisms.

\begin{hypothesis}[H-VRG2]
Phonon drag in a Gao--Khalaf-type moir\'e electron--phonon bath places 2D phase vortices in an overdamped regime $\hbar/\tau_v\gtrsim\pi\kappa$, converting a universal jump into a rounded crossover rather than a textbook thermal KT transition.
\end{hypothesis}

Two labeled mobility estimates are compared with the primary $\kappa(\alpha)$ at $T^*$. The phonon amplitude already present in the truncated GPE is $\hbar/\tau_v\sim g_{\mathrm{ph}}=0.1\Delta_0$. An Eliashberg-type OOM with $\lambda_{\mathrm{ep}}=1$ is $\hbar/\tau_v\sim 2\pi\lambda_{\mathrm{ep}}k_BT$. Along the primary branch both give $(\hbar/\tau_v)/(\pi\kappa)\sim 0.04$--$0.19<1$: the vortices are underdamped. A sharp stiffness jump could therefore survive if a KT surface were hit. H-VRG2 is not forced by these estimates; it remains a conjecture pending a microscopically computed vortex mobility.

\subsection{Observables, acceptance, and closures}
\label{sec:vrg_obs}

\begin{table}[htbp]
\centering
\caption{Fold versus vortex-RG observables, as tested by suite C1--C6. Acceptance of H-VRG1 on a given protocol requires the essential-singularity comparison on vortex-inclusive data over $\gtrsim 1.5$ decades in $t_A=(A_c-A)/A_c$, together with $\eta\to 1/4$ in $g(r)$.}
\label{tab:vrg_obs}
\small
\begin{tabular}{@{}llll@{}}
\toprule
Observable & Fold (established) & Vortex RG (conjectured) & Present suite \\
\midrule
$\Delta(A)$ & ends at $\Delta_c$; $\sqrt{t_A}$ & $\Delta\to 0$; $e^{-C/\sqrt{t_A}}$ & C4 bound at $T^*$; C5 on fold data \\
$\kappa$ & need not jump & universal jump~\eqref{eq:universal_jump} & C1 two-sector $\kappa(\alpha)$ \\
$g(r)$ & amplitude death first & algebraic $\to$ exponential & spin-wave at primary $K$ \\
vortex density & not required & proliferates above $A_c$ & C2 core; C3 plasma flag \\
mobility & --- & jump vs AHNS crossover & C6 underdamped at $T^*$ \\
specific heat & mean-field anomaly possible & no latent heat & program \\
\bottomrule
\end{tabular}
\end{table}

The vortex-RG sector is closed by the following statements. The Villain reduction from the \TQGL{} functional is an ansatz. The C1 identity~\eqref{eq:kappa} is derived given the healing-disk $n_0$ and the frozen geometric floor, both of which are ansatzes. Hypothesis H-VRG1 is a conjecture: if the drive path satisfies Theorem~\ref{thm:competing}, the gap closing takes the essential-singularity form~\eqref{eq:v9}. On the primary closure at $T^*$ the fold is reached first; conventional-only stiffness, and the primary closure at $T_{\mathrm{eff}}\sim 0.5$--$1$\,K, populate the KT side of the diagram as finite-$y$ plasmas. Chern--Simons remains truncated from the default scan unless Eq.~\eqref{eq:Kc_anyon} is restored as a labeled option. The locked amplitude fold is the mean-field backbone of the static theory and is independent of whether a given protocol realizes H-VRG1.
